%! Dividing Polynomials — Unit 3, Lesson 3.3 — Extra Credit
% UNITS 3+ PACKET-FIRST shape (LESSON_SHAPE.md §0a, 2026-11-14). GENERATED, GRADED, OPTIONAL.
% Last in the packet; the build starts it on a fresh sheet so it tears off. TWO PAGES, exactly,
% blank and key — 2pp is a ceiling. Five stretch items, harder than every ★★★:
%   1. COMBINE — long division (Round 1) with zero placeholders in dividend and divisor (Round 3)
%   2. MODEL — a riding-team feed bin: divide the volume by the height (Round 2), interpret R = 0
%      and check the volume at a given height (rubric)
%   3. REVERSE — rebuild P from Q and R; find k so the remainder is 0 (Rounds 2 + 3)
%   4. FIND THE MISTAKE — r = -4 used for x - 4 (the Round 2 flag), written response
%   5. SOL-style MULTIPLE CHOICE — the synthetic setup with placeholders (the crux)
% Every final answer is a \gradedblank, the line copied VERBATIM into the key (§9). Work blocks
% sit OUTSIDE the boxes. NO NAME LINE. Never circle/underline/bubble. Contexts differ from the
% notes' (study hall, the spring formal) — teach, then transfer. Points: 2+1+3+2+2+3+1 = 14.
% Assignment id: alg2-u03-l03-ec.
\documentclass[12pt]{article}
\usepackage{algebra2-article}
\usepackage{algebra2-boxes}
\usepackage{grading}          % math-shared: \gradedblank, QR codes, marks
\gradingassignment{alg2-u03-l03-ec}
\usepackage{algebra2-graded}
\gradingcourse{algebra_2}

\begin{document}

\pageheader{Unit 3, Lesson 3.3}{Extra Credit}

\vspace{0.10in}

\boxguard[8]
\begin{remindbox}
\small
\textbf{Extra credit --- optional.} Tear off these pages and turn them in with your homework,
the first class after two study halls. I will announce the due date in class and on TurtleNet.
{\small\itshape Write each final answer in its box. Only what is inside a box is graded.}\par
\end{remindbox}

\vspace{0.02in}

\begin{notesbox}{Stretch}
\small
\begin{enumerate}[leftmargin=1.9em, itemsep=7pt]
  \item \textbf{Two rounds at once.} Use long division, with a zero placeholder wherever a power
    is missing: $(x^4 + 3x^2 - 2x + 5) \div (x^2 - x + 1)$. Answer as $Q + R/D$.
    \begin{work}
      &\renewcommand{\arraystretch}{1.0}
      \begin{array}{r@{\;}r@{\;}r@{\;}r@{\;}r@{\;}r}
        & & & x^2 & +\,x & +\,3 \\ \cline{2-6}
        x^2-x+1\,\big) & x^4 & +\,0x^3 & +\,3x^2 & -\,2x & +\,5 \\
        & -(x^4 & -\,x^3 & +\,x^2) & & \\ \cline{2-4}
        & & x^3 & +\,2x^2 & -\,2x & \\
        & & -(x^3 & -\,x^2 & +\,x) & \\ \cline{3-5}
        & & & 3x^2 & -\,3x & +\,5 \\
        & & & -(3x^2 & -\,3x & +\,3) \\ \cline{4-6}
        & & & & & 2
      \end{array}
    \end{work}
    Answer: \inblank{\blank{6.4cm}}{\gradedblank[id=q1, type=short, points=2, answer={x^2+x+3+2/(x^2-x+1)}, check=expression, display={$x^2+x+3+2/(x^2-x+1)$}]{6.4cm-2pt}{12.5pt}}

  \item \textbf{Model it --- the riding team's feed bin.} The riding team stores feed in a bin
    shaped like a rectangular box. Its volume is $V(x) = x^3 + 7x^2 + 14x + 8$ cubic feet, and
    its height is $x + 2$ feet.
    \begin{enumerate}[label=(\alph*), leftmargin=1.6em, itemsep=6pt]
      \item Divide to find the area of the bin's base, as a polynomial.
        \begin{work}
          &\renewcommand{\arraystretch}{1.0}
          \begin{array}{r|rrrr}
            -2 & 1 & 7 & 14 & 8 \\
               &   & -2 & -10 & -8 \\ \cline{2-5}
               & 1 & 5 & 4 & 0
          \end{array}
          \qquad \text{base area} = \frac{V}{\text{height}} = x^2+5x+4
        \end{work}
        Base area: \inblank{\blank{3.6cm}}{\gradedblank[id=q2a, type=short, points=1, answer={x^2+5x+4}, check=expression]{3.6cm-2pt}{12.5pt}} square feet
      \item The remainder in (a) is $0$. Explain what that tells you about the bin. Then find the
        area of the base when the bin is $5$ feet tall, and use it to check the bin's volume.
        \rubricbox{\gradedblank[id=q2b, type=rubric, points=3, display={A remainder of 0 means $V(x) = (x+2)(x^2+5x+4)$ exactly: the volume is the height times the base area, as it must be for a box. Height 5 ft means $x=3$, so the base area is $9+15+4 = 28$ square feet, and $5 \cdot 28 = 140 = V(3) = 27+63+42+8$.}]{\linewidth}{2in}}
    \end{enumerate}
\end{enumerate}
\end{notesbox}

\newpage
\begin{notesbox}{Stretch, continued}
\small
\begin{enumerate}[leftmargin=1.9em, itemsep=10pt, start=3]
  \item \textbf{Run it backwards.}
    \begin{enumerate}[label=(\alph*), leftmargin=1.6em, itemsep=6pt]
      \item When $P(x)$ is divided by $x - 3$, the quotient is $x^2 + 2x - 1$ and the
        remainder is $4$. Find $P(x)$ in standard form.
        \begin{work}
          P(x) &= (x-3)(x^2+2x-1)+4 \qquad \text{($P = D\cdot Q + R$)} \\
          &= x^3+2x^2-x-3x^2-6x+3+4 = x^3-x^2-7x+7
        \end{work}
        $P(x) = $ \inblank{\blank{4.4cm}}{\gradedblank[id=q3a, type=short, points=2, answer={x^3-x^2-7x+7}, check=expression]{4.4cm-2pt}{12.5pt}}
      \item Find the number $k$ for which $(x^3 + kx^2 - 4) \div (x - 2)$ leaves a remainder of
        $0$. \textit{(Mind the missing term.)}
        \begin{work}
          &\renewcommand{\arraystretch}{1.2}
          \begin{array}{r|rrrr}
            2 & 1 & k & 0 & -4 \\
              &   & 2 & 2k+4 & 4k+8 \\ \cline{2-5}
              & 1 & k+2 & 2k+4 & 4k+4
          \end{array}
          \qquad 4k+4=0 \;\Rightarrow\; k=-1
        \end{work}
        $k = $ \inblank{\blank{1.8cm}}{\gradedblank[id=q3b, type=short, points=2, answer={-1}, check=numeric]{1.8cm-2pt}{12.5pt}}
    \end{enumerate}

  \item \textbf{Find the mistake.} Student A divided $(x^3 - 6x^2 + 32) \div (x - 4)$:
    \begin{center}
    \renewcommand{\arraystretch}{1.3}
    \begin{tabular}{r|rrrr}
      $-4$ & $1$ & $-6$ & $0$ & $32$ \\
           &     & $-4$ & $40$ & $-160$ \\ \cline{2-5}
           & $1$ & $-10$ & $40$ & $-128$
    \end{tabular}
    \qquad
    \begin{minipage}{0.42\linewidth}
    ``So the answer is $x^2 - 10x + 40 - \dfrac{128}{x-4}$.''
    \end{minipage}
    \end{center}
    Explain Student A's mistake, then give the correct quotient and remainder.
    \rubricbox{\gradedblank[id=q4, type=rubric, points=3, display={She used $r=-4$, but for $x-4$ the number is $r=4$ (the value that makes $x-4$ zero). Correct: $4 \mid 1,\ -6,\ 0,\ 32$ gives $1,\ -2,\ -8 \mid 0$, so the quotient is $x^2-2x-8$ and the remainder is $0$ --- it divides evenly. Check: $(x-4)(x^2-2x-8) = x^3-6x^2+32$.}]{\linewidth}{2in}}

  \item \textbf{Multiple choice (SOL-style).} Which shows the correct first row for the
    synthetic division $(3x^4 - 2x + 7) \div (x + 5)$? Write the letter of the single best
    answer in its Letter box.
    \begin{enumerate}[label=(\Alph*), leftmargin=2.0em, itemsep=3pt]
      \item $\phantom{-}5 \;\big|\; 3 \quad 0 \quad 0 \quad -2 \quad 7$
      \item $-5 \;\big|\; 3 \quad -2 \quad 7$
      \item $-5 \;\big|\; 3 \quad 0 \quad 0 \quad -2 \quad 7$
      \item $-5 \;\big|\; 3 \quad 0 \quad -2 \quad 7$
    \end{enumerate}
    \letterbox{\gradedblank[id=q5, type=choice, points=1, answer={C}]{0.4in}{0.4in}}
\end{enumerate}
\end{notesbox}

\end{document}
